\documentclass[11pt,letterpaper]{article}
\usepackage[utf8]{inputenc}
\usepackage[spanish,es-noshorthands]{babel}
\usepackage{amsmath,amssymb}
\usepackage{graphicx}
\usepackage{xcolor}
\usepackage{tikz}
\usetikzlibrary{calc,arrows.meta,decorations.markings}
\usepackage[top=1.8cm, bottom=2.2cm, left=2.5cm, right=2.5cm]{geometry}
\pagestyle{plain}

% simbolos: punto (sale de la pagina) y cruz (entra a la pagina)
\newcommand{\punto}[2][black]{%
  \draw[#1,line width=0.7pt,fill=white] (#2) circle (0.16);
  \fill[#1] (#2) circle (0.045);}
\newcommand{\cruz}[2][black]{%
  \draw[#1,line width=0.7pt,fill=white] (#2) circle (0.16);
  \draw[#1,line width=0.7pt] ($(#2)+(-0.11,-0.11)$)--($(#2)+(0.11,0.11)$)
        ($(#2)+(-0.11,0.11)$)--($(#2)+(0.11,-0.11)$);}
% sistema de referencia: x derecha, y arriba, z saliendo de la pagina
\newcommand{\ejesxyz}[1]{%
  \begin{scope}[shift={(#1)}]
    \draw[->,line width=0.7pt] (0,0)--(0.75,0) node[right,font=\footnotesize]{$x$};
    \draw[->,line width=0.7pt] (0,0)--(0,0.75) node[above,font=\footnotesize]{$y$};
    \draw[line width=0.6pt,fill=white] (0,0) circle (0.1); \fill (0,0) circle (0.03);
    \node[below left=-1pt,font=\footnotesize] at (0,0){$z$};
  \end{scope}}

\begin{document}

\begin{center}
\includegraphics[height=1.9cm]{logo ubb.png}\\[0.25cm]
\textbf{F\'isica II (2026-2)}\\
\textbf{Tarea 4}\emph{: Ley de Biot-Savart, Ley de Amp\`ere y Ley de Faraday-Lenz}\\
\textit{Profesor: Arturo Fern\'andez P\'erez}\\
\textbf{Plazo de Entrega: Mi\'ercoles 18 de noviembre a las 23:59 hrs.}
\end{center}

\begin{enumerate}
\renewcommand{\labelenumii}{(\alph{enumii})}

%==================================================================
% PROBLEMA 1
%==================================================================
\item Tres cables rectos muy largos y paralelos al eje $y$ se encuentran en el plano $xy$,
ubicados en $x_1 = 0$, $x_2 = 10$ [cm] y $x_3 = 25$ [cm], como muestra la Fig.\ 1. Los cables
conducen corrientes $I_1 = 2$ [A] en la direcci\'on $+y$, $I_2 = 5$ [A] en la direcci\'on $-y$ e
$I_3 = 3$ [A] en la direcci\'on $-y$. Considerando el sistema de referencia de la figura
(eje $z^{+}$ saliendo de la p\'agina), determine:

\begin{enumerate}
\item Usando la ley de Amp\`ere, la magnitud del campo magn\'etico que cada cable produce en la
posici\'on de los otros dos.

\item Usando la regla de la mano derecha, la direcci\'on de cada uno de esos campos, expres\'andola
con la notaci\'on de puntos y cruces y como vector ($\pm\hat{k}$).

\item Usando $\vec{F} = I\,\vec{L}\times\vec{B}$, el vector fuerza magn\'etica por unidad de longitud
sobre cada cable.
\end{enumerate}

\begin{center}
\begin{tikzpicture}[scale=0.9, >=Stealth, line width=0.9pt]
    \foreach \x/\n in {0/1, 2/2, 5/3}{
        \draw[orange!75!black,line width=2.2pt] (\x,-2.0)--(\x,2.0);
        \node[font=\small] at (\x,-2.35){$(\n)$};}
    % corrientes
    \draw[->,violet,line width=1.3pt] (-0.3,-0.4)--(-0.3,0.6);
    \node[violet,left,font=\small] at (-0.3,0.1){$I_1$};
    \draw[->,violet,line width=1.3pt] (1.7,0.6)--(1.7,-0.4);
    \node[violet,left,font=\small] at (1.7,0.1){$I_2$};
    \draw[->,violet,line width=1.3pt] (4.7,0.6)--(4.7,-0.4);
    \node[violet,left,font=\small] at (4.7,0.1){$I_3$};
    % distancias
    \draw[<->,gray] (0,1.6)--(2,1.6);  \node[above,font=\footnotesize,gray] at (1,1.6){$10$ cm};
    \draw[<->,gray] (2,1.6)--(5,1.6);  \node[above,font=\footnotesize,gray] at (3.5,1.6){$15$ cm};
    % sistema de referencia
    \ejesxyz{6.2,-1.8}
\end{tikzpicture}\\[2pt]
Figura 1
\end{center}

%==================================================================
% PROBLEMA 2
%==================================================================
\item Dos alambres con forma de arco circular, centrados en el origen $O$ y contenidos en el plano
$xy$, conducen corrientes en sentidos opuestos, como muestra la Fig.\ 2. El arco interior tiene radio
$R_1 = 6$ [cm], subtiende un \'angulo de $30^\circ$ y conduce una corriente $I_1 = 4$ [A] en sentido
antihorario. El arco exterior tiene radio $R_2 = 10$ [cm], subtiende un \'angulo de $120^\circ$ y
conduce una corriente $I_2 = 3$ [A] en sentido horario. Considerando s\'olo el campo producido por
los arcos, determine:

\smallskip
\noindent
\begin{minipage}[c]{0.50\linewidth}
\begin{enumerate}
\item Usando la ley de Biot-Savart, el campo magn\'etico (magnitud y direcci\'on) que produce cada
arco en el origen.

\item El campo magn\'etico neto en el origen, indicando si apunta en la direcci\'on $+z$ o $-z$.
\end{enumerate}
\end{minipage}%
\hfill
\begin{minipage}[c]{0.46\linewidth}
\centering
\begin{tikzpicture}[scale=0.72, >=Stealth, line width=0.9pt]
    % ejes
    \draw[->,line width=0.7pt] (-3.9,0)--(3.2,0) node[right,font=\small]{$x$};
    \draw[->,line width=0.7pt] (0,-3.3)--(0,3.4) node[above,font=\small]{$y$};
    % radios auxiliares (no son alambres)
    \draw[gray,dashed,line width=0.6pt] (0,0)--(30:1.8);
    \draw[gray,dashed,line width=0.6pt] (0,0)--(120:3.0);
    \draw[gray,dashed,line width=0.6pt] (0,0)--(240:3.0);
    % arco 1: R1, de 0 a 30 grados, antihorario
    \draw[orange!75!black,line width=2.2pt,postaction={decorate},
          decoration={markings, mark=at position 0.65 with {\arrow[violet,scale=1.1]{>}}}]
          (0:1.8) arc[start angle=0, end angle=30, radius=1.8];
    \node[violet,font=\small] at (18:2.3){$I_1$};
    \node[font=\footnotesize] at (1.25,-0.3){$R_1$};
    % arco 2: R2, de 240 a 120 grados, horario
    \draw[orange!75!black,line width=2.2pt,postaction={decorate},
          decoration={markings, mark=at position 0.5 with {\arrow[violet,scale=1.1]{>}}}]
          (240:3.0) arc[start angle=240, end angle=120, radius=3.0];
    \node[violet,font=\small] at (172:3.45){$I_2$};
    \node[font=\footnotesize,left] at (120:2.0){$R_2$};
    % angulos
    \draw[line width=0.6pt] (0:0.7) arc[start angle=0, end angle=30, radius=0.7];
    \node[font=\scriptsize] at (15:1.08){$30^\circ$};
    \draw[line width=0.6pt] (120:0.7) arc[start angle=120, end angle=240, radius=0.7];
    \node[font=\scriptsize] at (200:1.2){$120^\circ$};
    % origen
    \fill (0,0) circle (2pt);
    \node[below left,font=\small] at (0,0){$O$};
    % z saliendo
    \draw[line width=0.6pt,fill=white] (2.6,2.8) circle (0.1); \fill (2.6,2.8) circle (0.03);
    \node[right,font=\footnotesize] at (2.7,2.8){$z$};
\end{tikzpicture}\\[2pt]
Figura 2
\end{minipage}

\pagebreak
%==================================================================
% PROBLEMA 3
%==================================================================
\item Dos cables rectos muy largos y perpendiculares al plano de la hoja est\'an separados una
distancia $d = 5$ [cm], como muestra la Fig.\ 3. El cable 1 conduce una corriente $I_1 = 6$ [A]
\textbf{entrando} al plano de la hoja y el cable 2 conduce una corriente $I_2 = 8$ [A]
\textbf{saliendo} del plano de la hoja. El punto $P$ se encuentra a $3$ [cm] del cable 1 y a $4$
[cm] del cable 2. Determine:

\begin{enumerate}
\item Usando la ley de Amp\`ere, la magnitud del campo magn\'etico que produce cada cable en el
punto $P$.

\item Usando la regla de la mano derecha, la direcci\'on de cada campo en $P$. Dib\'ujelos en la
figura y exprese cada uno en t\'erminos de $\hat{\imath}$ y $\hat{\jmath}$.

\item El campo magn\'etico neto en $P$, su magnitud y el \'angulo que forma con el eje $x$.
\end{enumerate}

\begin{center}
\begin{tikzpicture}[scale=1.0, >=Stealth, line width=0.9pt]
    % lados del triangulo
    \draw[gray,dashed,line width=0.7pt] (0,0)--(1.8,2.4)--(5,0);
    \draw[<->,gray,line width=0.6pt] (0,-0.5)--(5,-0.5);
    \node[below,font=\footnotesize,gray] at (2.5,-0.5){$d = 5$ cm};
    \node[font=\footnotesize,gray] at (0.55,1.35){$3$ cm};
    \node[font=\footnotesize,gray] at (3.75,1.45){$4$ cm};
    % angulo recto en P
    \draw[gray,line width=0.5pt] ($(1.8,2.4)!0.3cm!(0,0)$) -- ++(0.8*0.3,-0.6*0.3) -- ($(1.8,2.4)!0.3cm!(5,0)$);
    % cables
    \cruz[orange!75!black]{0,0}
    \punto[orange!75!black]{5,0}
    \node[below left,font=\small] at (-0.1,-0.05){$I_1$};
    \node[below right,font=\small] at (5.1,-0.05){$I_2$};
    \node[above left,font=\footnotesize] at (-0.1,0.1){$(1)$};
    \node[above right,font=\footnotesize] at (5.1,0.1){$(2)$};
    % punto P
    \fill (1.8,2.4) circle (2pt);
    \node[above,font=\small] at (1.8,2.45){$P$};
    % sistema de referencia
    \ejesxyz{6.3,1.6}
\end{tikzpicture}\\[2pt]
Figura 3
\end{center}

%==================================================================
% PROBLEMA 4
%==================================================================
\item Un solenoide muy largo tiene $n = 4000$ [vueltas/m], radio $R = 3$ [cm] y conduce una
corriente variable $I(t) = 3\sqrt{t}$ [A], con $t$ en segundos. Dos bobinas son conc\'entricas y
coaxiales con el solenoide (Fig.\ 4): una bobina \textbf{interior} de $N_1 = 150$ vueltas, radio
$r_1 = 2$ [cm] y resistencia $R_1 = 0{,}5$ [$\Omega$], y una bobina \textbf{exterior} de
$N_2 = 100$ vueltas, radio $r_2 = 5$ [cm] y resistencia $R_2 = 1{,}2$ [$\Omega$]. Determine:

\begin{enumerate}
\item El campo magn\'etico en el interior del solenoide en $t = 4$ [s].

\item El flujo magn\'etico a trav\'es de una vuelta de cada bobina, en funci\'on del tiempo.

\item La fem inducida y la corriente inducida en cada bobina en $t = 4$ [s].

\item Usando la ley de Lenz, si la corriente inducida en cada bobina circula en el mismo sentido
o en sentido opuesto a la corriente del solenoide. Justifique.
\end{enumerate}

\begin{center}
\begin{tikzpicture}[scale=0.9, >=Stealth, line width=0.9pt]
    %----------------- vista lateral -----------------
    \begin{scope}
        % bobina exterior (mitad trasera)
        \foreach \x in {-0.45,0,0.45}{
            \draw[green!50!black,line width=1.1pt] (\x,1.7) arc[start angle=90, end angle=-90,
                  x radius=0.22, y radius=1.7];}
        % solenoide (espiras delgadas)
        \foreach \x in {-2.55,-2.25,...,2.6}{
            \draw[blue!70!black,line width=0.6pt] (\x,0) ellipse[x radius=0.1, y radius=1.0];}
        % bobina interior
        \foreach \x in {-0.3,0,0.3}{
            \draw[red!80!black,line width=1.3pt] (\x+0.02,0) ellipse[x radius=0.08, y radius=0.67];}
        % bobina exterior (mitad delantera)
        \foreach \x in {-0.45,0,0.45}{
            \draw[green!50!black,line width=1.1pt] (\x,-1.7) arc[start angle=-90, end angle=-270,
                  x radius=0.22, y radius=1.7];}
        % eje
        \draw[gray,dash dot,line width=0.5pt] (-3.1,0)--(3.2,0);
        % corriente del solenoide
        \draw[->,violet,line width=1.2pt] (-3.3,-1.0)--(-2.7,-1.0);
        \node[violet,left,font=\small] at (-3.3,-1.0){$I(t)$};
        \node[blue!70!black,font=\footnotesize] at (-1.6,1.35){solenoide};
        \draw[red!80!black,line width=0.5pt] (0.35,-0.45)--(1.1,-1.95);
        \node[red!80!black,font=\footnotesize,below] at (1.1,-1.95){$N_1$};
        \node[green!50!black,font=\footnotesize] at (1.15,1.85){$N_2$};
        \node[font=\footnotesize] at (-1.2,-2.25){Vista lateral};
    \end{scope}
    %----------------- vista frontal -----------------
    \begin{scope}[shift={(6.6,0)}]
        \draw[green!50!black,line width=1.3pt] (0,0) circle (1.7);
        \draw[blue!70!black,line width=1.3pt,postaction={decorate},
              decoration={markings, mark=at position 0.1 with {\arrow[violet,scale=1.2]{>}}}]
              (0,0) circle (1.0);
        \draw[red!80!black,line width=1.3pt] (0,0) circle (0.67);
        \fill (0,0) circle (1.5pt);
        \draw[->,line width=0.6pt] (0,0)--(-60:0.67);  \node[font=\scriptsize] at (-35:0.45){$r_1$};
        \draw[->,line width=0.6pt] (0,0)--(160:1.0);   \node[font=\scriptsize,above] at (160:0.62){$R$};
        \draw[->,line width=0.6pt] (0,0)--(215:1.7);   \node[font=\scriptsize] at (225:1.3){$r_2$};
        \node[violet,font=\small] at (52:1.32){$I$};
        \node[font=\footnotesize,align=center] at (0,-2.25){Vista frontal\\(corriente antihoraria)};
    \end{scope}
\end{tikzpicture}\\[2pt]
Figura 4
\end{center}

\pagebreak
%==================================================================
% PROBLEMA 5
%==================================================================
\item Una bobina rectangular de $N = 50$ vueltas y lados $a = 20$ [cm] y $b = 15$ [cm] gira con
rapidez angular constante $\omega = 30$ [rad/s] alrededor de un eje perpendicular a un campo
magn\'etico uniforme $\vec{B}$ de magnitud desconocida, como muestra la Fig.\ 5. La resistencia
total de la bobina es $R = 12$ [$\Omega$] y la corriente inducida m\'axima es
$I_{\text{m\'ax}} = 0{,}45$ [A]. Determine:

\begin{enumerate}
\item La fem inducida m\'axima en la bobina.

\item Usando la ley de Faraday, la magnitud del campo magn\'etico $B$.

\item Las expresiones del flujo magn\'etico $\Phi_B(t)$ a trav\'es de una vuelta y de la fem
inducida $\varepsilon(t)$, si en $t = 0$ el plano de la bobina es perpendicular a $\vec{B}$.
\end{enumerate}

\begin{center}
\begin{tikzpicture}[scale=0.9, >=Stealth, line width=0.9pt]
    % lineas de campo uniforme
    \foreach \y in {-1.8,-0.9,0,0.9,1.8}{
        \draw[->,blue!70!black,line width=0.8pt] (-3.2,\y)--(3.4,\y);}
    \node[blue!70!black,right,font=\small] at (3.4,1.8){$\vec{B}$};
    % eje de rotacion (parte trasera)
    \draw[gray,dash dot,line width=0.7pt] (0,-2.5)--(0,2.6);
    % bobina (rectangulo girado en torno al eje vertical)
    \draw[orange!75!black,line width=2.2pt,fill=orange!12,fill opacity=0.85]
          (-1.1,-1.25)--(1.1,-0.95)--(1.1,1.45)--(-1.1,1.15)--cycle;
    \draw[gray,dash dot,line width=0.7pt] (0,-1.1)--(0,1.3);
    % cotas
    \node[font=\small] at (1.4,0.25){$a$};
    \node[font=\small] at (0.2,-1.45){$b$};
    % rapidez angular
    \draw[->,line width=1pt] (0.55,2.3) arc[start angle=-30, end angle=210, x radius=0.6, y radius=0.2];
    \node[font=\small] at (1.05,2.55){$\omega$};
\end{tikzpicture}\\[2pt]
Figura 5
\end{center}

\end{enumerate}

\end{document}
